%% Draft paper for SAE-SMO-SPLADE experiments.
%% This file intentionally contains both write-up text and TODO placeholders so
%% future runs can be dropped into the tables without rethinking the structure.
\documentclass[
% twocolumn,
% hf,
]{ceurart}

\sloppy

\usepackage{booktabs}
\usepackage{listings}
\usepackage{amsmath}
\usepackage{array}
\usepackage{multirow}
\lstset{breaklines=true,basicstyle=\ttfamily\small}

\begin{document}

\copyrightyear{2026}
\copyrightclause{Copyright for this paper by its authors.
  Use permitted under Creative Commons License Attribution 4.0
  International (CC BY 4.0).}

\conference{TODO: Workshop / venue name, TODO date, TODO location}

\title{Lightweight Query Encoders for Frozen Sparse Document Models}

\author[1]{TODO Author One}[%
email=todo@example.com,
]
\author[1]{TODO Author Two}[%
email=todo@example.com,
]
\address[1]{TODO Affiliation}

\begin{abstract}
Learned sparse retrievers such as SPLADE provide strong lexical matching by
expanding queries and documents into a shared vocabulary space, but standard
usage applies the same relatively expensive encoder to both queries and
documents. This work studies an asymmetric alternative: keep a strong document
encoder frozen and offline, while training a cheaper query encoder whose sparse
vectors remain directly dot-product compatible with the document index. We
report a sequence of model-surgery experiments: vocabulary transplant between
different masked-language-model vocabularies, shallow query encoders obtained by
truncating a full SPLADE document model, ALBERT-style factorization of the
shared embedding/MLM-head matrix, and extensions to Lion-SP-1B. The most stable
SPLADE-v3 result uses a 3-layer spaced query body with an SVD-initialized
factorized lexical head kept frozen during training; on full MS MARCO dev it
achieves NDCG@10 0.4438 and MRR@10 0.3787, compared with 0.4657 and 0.3989 for
the full SPLADE-v3 doc-only ceiling. This keeps about 95\% of the symmetric
retrieval quality while substantially reducing query-side parameters. The
ablation suite shows that cross-vocabulary vocabulary transplant can work, but
its embedding initialization is load-bearing; random initialization, frozen
document heads without transplant, and shared-vocabulary BERT baselines do not
match the transplanted model. We also document failure modes: cross-encoder
MarginMSE distillation can produce good retrieval while yielding unstable sparse
support, and Lion-scale tokenizers make cross-vocabulary transfer especially
fragile. The paper is a draft template; TODO rows mark active and future
experiments.
\end{abstract}

\begin{keywords}
  sparse retrieval \sep
  SPLADE \sep
  asymmetric retrieval \sep
  model compression \sep
  vocabulary transfer \sep
  knowledge distillation
\end{keywords}

\maketitle

\section{Paper Plan}

This section is a working plan for finishing the paper. It can be removed or
moved to a private notes file before submission.

\paragraph{Primary story.}
The cleanest story is not that every experiment succeeded. The main story is
that dot-product-compatible asymmetric sparse retrieval is feasible, but the
route matters. Cross-vocabulary transplant is possible and instructive, yet
brittle. The strongest practical path is to start from the same document model,
truncate the query body, and compress the lexical matrix while preserving the
document vocabulary.

\paragraph{Recommended contribution framing.}
\begin{enumerate}
  \item Define the asymmetric sparse retrieval setting: frozen expensive
  document encoder, cheap query encoder, same sparse vocabulary, no projection
  at retrieval time.
  \item Present two families of query-side model surgery:
  cross-vocabulary transplant and same-checkpoint shallow truncation.
  \item Show that transplant initialization quality is the binding constraint
  in cross-vocabulary runs.
  \item Show that factorized shallow SPLADE preserves most of the full
  SPLADE-v3 MS MARCO quality.
  \item Document negative results as engineering and scientific evidence:
  random embeddings, frozen document heads, shared-vocabulary BERT, and
  cross-encoder MarginMSE are not reliable replacements for the stable recipe.
\end{enumerate}

\paragraph{Experiments still needed before submission.}
\begin{enumerate}
  \item Full MS MARCO dev for the best Lion query checkpoint. This requires a
  Lion-specific 128K-vocabulary document index.
  \item Repeat the best SPLADE-v3 factorized result with the corrected
  optimizer-step learning-rate schedule, since gradient accumulation currently
  makes cosine decay much slower than intended in Lion runs.
  \item Add latency and parameter-count tables for all reported checkpoints.
  \item Add a seed or repeatability check for the strongest NanoBEIR-only
  claims, or clearly mark them as exploratory.
  \item Decide whether the paper title should emphasize ``model surgery'',
  ``asymmetric sparse retrieval'', or ``query-side compression''.
\end{enumerate}

\section{Introduction}

Neural sparse retrievers bridge lexical retrieval and neural semantic matching
by producing high-dimensional sparse vectors in a fixed vocabulary space.
SPLADE-style models \cite{Formal2021SPLADE,Formal2022SPLADEv2} use masked
language-model logits and max pooling to expand text into weighted lexical
features. This makes retrieval compatible with inverted indexes while achieving
strong ranking performance.

However, standard SPLADE deployment uses the same encoder for queries and
documents. Document encoding can be performed offline, but query encoding
remains an online cost. The serving problem becomes sharper for stronger sparse
document models, including SPLADE-v3 \cite{Lassance2024SPLADEv3} and recent
LLM-based sparse retrievers such as Lion-SP. The core question of this work is:
can we keep the strong document encoder fixed and train a much cheaper query
encoder that emits vectors in the same sparse space?

This setting has a strict compatibility requirement. If the query and document
encoders do not share the same output vocabulary, their sparse vectors cannot be
directly dot-producted. We therefore study methods that preserve or create a
shared sparse vocabulary:

\begin{itemize}
  \item \textbf{Vocabulary transplant}: move the document encoder's tokenizer
  and lexical matrix onto a smaller query backbone using token-level embedding
  interpolation.
  \item \textbf{Shallow same-checkpoint query encoders}: initialize the query
  model from the document model, keep the same tokenizer and lexical head, and
  remove most transformer layers.
  \item \textbf{Factorized lexical matrices}: reduce the parameter floor caused
  by the tied input embedding and MLM head using an ALBERT-style low-rank
  factorization \cite{Lan2020ALBERT}.
\end{itemize}

The empirical arc is important. The first transplant approach revealed that
embedding initialization is crucial: correct vocabulary size alone is not
enough. The later shallow-query approach avoids that brittleness by construction
and currently gives the strongest practical result.

\section{Background and Related Work}

\subsection{Learned Sparse Retrieval}

SPLADE \cite{Formal2021SPLADE,Formal2022SPLADEv2} represents text by applying a
sparse transformation to masked-language-model logits. Given token-level logits
$z_{t,v}$ for token position $t$ and vocabulary term $v$, a typical SPLADE
encoding is:
\begin{equation}
  s_v(x) = \log(1 + \max_t \operatorname{ReLU}(z_{t,v})).
\end{equation}
This produces non-negative sparse lexical vectors that can be indexed with
inverted-index infrastructure. SPLADE-v3 \cite{Lassance2024SPLADEv3} improves
training and distillation recipes for SPLADE models.

Other learned sparse retrieval systems include TILDE/TILDEv2
\cite{Zhuang2021TILDE}, uniCOIL \cite{Lin2021UniCOIL}, DeepImpact
\cite{Mallia2021DeepImpact}, and lexical bottleneck systems such as LexMAE
\cite{Shen2023LexMAE}. These methods share the goal of producing efficient
retrieval-time representations, but our focus is specifically asymmetric
query-side compression while keeping a strong frozen document encoder fixed.

\subsection{Late Interaction and Distillation}

ColBERT and ColBERTv2 \cite{Santhanam2022ColBERTv2} show that high-quality
retrieval can be achieved by separating query and document representations and
using efficient late interaction. SPLADE training recipes also rely heavily on
distillation from cross-encoders and rerankers. We attempted a traditional
MarginMSE-style cross-encoder distillation route, but found it under-constrained
for the asymmetric sparse setting: margins can be matched by many different
query-vector supports.

\subsection{Vocabulary and Embedding Transfer}

Vocabulary mismatch is a central obstacle in cross-model sparse retrieval.
WECHSEL \cite{Minixhofer2022WECHSEL} and related vocabulary-transfer methods
show that careful embedding initialization can transfer language models across
tokenizers. Our vocabulary-transplant experiments use the same broad idea:
shared tokens are copied directly, while donor-only tokens are initialized by
nearest-neighbor interpolation from shared-token geometry. This is not enough
for all backbones, but it is far better than random initialization.

\subsection{Embedding Factorization}

ALBERT \cite{Lan2020ALBERT} reduces parameters by factorizing the vocabulary
embedding matrix into a smaller lexical embedding and an up-projection to hidden
size. Our factorized SPLADE query models adapt this idea to a tied
input/output lexical matrix. The query encoder maps token ids through
$A \in \mathbb{R}^{|V| \times r}$ and $B \in \mathbb{R}^{r \times h}$, and the
output head maps hidden states back through $B^\top A^\top$. This preserves the
same output vocabulary while reducing the lexical parameter floor.

\section{Problem Setup}

Let $D_\theta$ be a strong frozen document encoder and $Q_\phi$ a cheaper query
encoder. Both must output sparse vectors in the same vocabulary:
\begin{equation}
  D_\theta(d), Q_\phi(q) \in \mathbb{R}^{|V|}_{\ge 0}.
\end{equation}
Documents are encoded offline and stored in a sparse index. At retrieval time,
the query encoder produces $Q_\phi(q)$ and documents are ranked by dot product:
\begin{equation}
  \operatorname{score}(q,d) = Q_\phi(q)^\top D_\theta(d).
\end{equation}

The objective is to reduce query-side compute and parameter count while
preserving the retrieval quality of a symmetric or doc-only ceiling in which
the strong document model is also used to encode queries.

\section{Methods}

\subsection{Vocabulary Transplant}

The vocabulary-transplant approach starts with a small query masked-language
model and a stronger SPLADE document model with a different tokenizer. The goal
is to replace the query model's vocabulary with the document model's vocabulary
without adding a projection layer.

For each donor vocabulary token $v$:
\begin{enumerate}
  \item If $v$ exists in both vocabularies, copy the corresponding query-model
  embedding.
  \item Otherwise, find the $k$ nearest shared tokens to $v$ in the donor
  embedding space.
  \item Interpolate the query embeddings of those shared neighbors using
  inverse-distance weights.
  \item Resize the query embedding and output head so the model emits logits in
  the donor vocabulary.
\end{enumerate}

In code, the interpolation is:
\begin{equation}
  e_Q(v) = \sum_{i \in \mathcal{N}_k(v)}
  \frac{(1 - \cos(e_D(v), e_D(i)) + \epsilon)^{-1}}
       {\sum_j (1 - \cos(e_D(v), e_D(j)) + \epsilon)^{-1}}
  e_Q(i).
\end{equation}

This procedure makes the transplanted query model dot-product compatible with
the frozen document index. The experiments show that this initialization is
load-bearing: random embeddings with the correct vocabulary size fail.

\subsection{Shallow Same-Checkpoint Query Encoders}

The more stable approach avoids cross-vocabulary transfer. We load the same
checkpoint for both query and document encoders, then remove most query-side
transformer layers while keeping the tokenizer, embedding table, and output
head. For SPLADE-v3, this gives a shallow BERT/SPLADE query encoder in the same
30,522-dimensional vocabulary as the full document encoder. For Lion-SP-1B, the
same idea gives a shallow Llama-style query encoder in a 128,256-dimensional
vocabulary.

Training uses a two-phase schedule:
\begin{enumerate}
  \item \textbf{Warmup}: freeze embeddings and the lexical head; train only the
  retained body layers.
  \item \textbf{Fine-tuning}: either keep the lexical factors/head frozen, or
  unfreeze them with a much smaller learning rate.
\end{enumerate}

This schedule was necessary in practice. Immediately training the lexical
matrix can destabilize sparse support; freezing the lexical basis lets the body
learn how to drive the original SPLADE vocabulary first.

\subsection{Factorized Lexical Head}

For BERT-sized SPLADE, the tied lexical matrix has
$30{,}522 \times 768 \approx 23.4$M parameters. For Lion-SP-1B, the tied Llama
lexical matrix has $128{,}256 \times 2048 \approx 262.7$M parameters. Removing
transformer layers therefore does not fully solve query-side parameter cost.

We replace the tied lexical matrix $W \in \mathbb{R}^{|V| \times h}$ with two
factors:
\begin{equation}
  W \approx A B,\quad A \in \mathbb{R}^{|V| \times r},\quad
  B \in \mathbb{R}^{r \times h}.
\end{equation}
The input embedding path uses $A B$, while the output head uses $B^\top A^\top$.
The factors are initialized with SVD from the original lexical matrix. The most
stable SPLADE-v3 factorized run keeps these factors frozen throughout training.

\subsection{Losses}

The main shallow-alignment loss matches the full document model's SPLADE query
vector:
\begin{equation}
  \mathcal{L}_{\mathrm{vec}} =
  \frac{1}{B}\sum_{i=1}^B \|Q_\phi(q_i) - D_\theta(q_i)\|_2^2.
\end{equation}
We use a straight-through estimator through ReLU so gradients are not completely
blocked by sparse thresholding early in training.

Sparsity is encouraged with query L1 mass:
\begin{equation}
  \mathcal{L}_{q} = \lambda_q \frac{1}{B}\sum_i \|Q_\phi(q_i)\|_1.
\end{equation}

The non-factorized SPLADE shallow run can also use a contrastive term over
Tevatron positives and hard negatives:
\begin{equation}
  \mathcal{L}_{\mathrm{ce}} =
  -\log
  \frac{\exp(Q_\phi(q_i)^\top D_\theta(d_i^+) / \tau)}
       {\sum_{j=1}^{n}\exp(Q_\phi(q_i)^\top D_\theta(d_{ij}) / \tau)}.
\end{equation}
This is expensive because it requires frozen document encodes during training,
so the strongest factorized SPLADE-v3 run disables it for cheaper iteration.

\section{Experimental Setup}

\subsection{Datasets and Metrics}

Training uses MS MARCO passage data \cite{Bajaj2016MSMARCO} through Tevatron
query-positive-hard-negative triples. Fast development uses NanoBEIR-style small
evaluation subsets of MS MARCO and NFCorpus. Full evaluation uses MS MARCO dev
with approximately 6,980 queries and the 8.8M-passage corpus. We report NDCG@10
and MRR@10.

NanoBEIR numbers are useful for iteration but noisy because each dataset has
about 50 queries. Full MS MARCO numbers should be treated as the main evidence
when available.

\subsection{Model Families}

\begin{table}[t]
\caption{Main model families studied in this draft. ``Status'' is intentionally
pragmatic: it describes whether the method is currently worth using.}
\label{tab:model-families}
\centering
\begin{tabular}{p{0.26\linewidth}p{0.48\linewidth}p{0.16\linewidth}}
\toprule
Family & Description & Status \\
\midrule
SAE-SPLADE & Top-$k$ sparse autoencoder over a small backbone, then SPLADE
fine-tuning. & Baseline \\
Vocabulary transplant & Small query backbone resized to document vocabulary
with kNN embedding interpolation. & Instructive, brittle \\
Projected query & Small backbone projected into frozen SPLADE head. & Negative \\
Direct shared-vocab BERT & BERT-family query model already sharing
SPLADE-v3 vocabulary. & Negative \\
Shallow SPLADE & Same checkpoint as doc model, but fewer query layers. & Stable \\
Factorized shallow SPLADE & Shallow SPLADE plus ALBERT-style lexical factors. & Strongest SPLADE-v3 result \\
Lion variants & Same ideas with Lion-SP-1B and 128K vocabulary. & Active \\
\bottomrule
\end{tabular}
\end{table}

\section{Results}

\subsection{Full MS MARCO Results}

\begin{table}[t]
\caption{Full MS MARCO dev results with 8.8M-passage indexes. These are the
strongest completed full-index results so far.}
\label{tab:full-msmarco}
\centering
\begin{tabular}{lrrp{0.35\linewidth}}
\toprule
Configuration & NDCG@10 & MRR@10 & Notes \\
\midrule
\texttt{naver/splade-v3} doc-only ceiling
  & 0.4657 & 0.3989 & Full SPLADE-v3 encodes queries and documents. \\
\texttt{splade\_shallow\_factorized\_spaced\_align}
  & 0.4438 & 0.3787 & 3-layer query body from layers [0, 6, 11]; frozen SVD lexical factors/head. \\
\texttt{lion\_shallow\_factorized\_align}
  & TODO & TODO & Requires Lion 128K-vocabulary full MS MARCO index. \\
\texttt{lion\_shallow\_factorized\_spaced\_align}
  & TODO & TODO & Active/future full-index eval after best checkpoint is chosen. \\
\bottomrule
\end{tabular}
\end{table}

The factorized spaced SPLADE-v3 query keeps approximately 95.3\% of the
SPLADE-v3 doc-only NDCG@10 and 94.9\% of its MRR@10. The absolute gaps are
0.0219 NDCG@10 and 0.0202 MRR@10.

\subsection{NanoBEIR Development Results}

\begin{table}[t]
\caption{Representative NanoBEIR development results. These numbers are useful
for comparing recipes during iteration, but should not be treated as final
claims without full MS MARCO evaluation or repeat runs.}
\label{tab:nanobeir}
\centering
\begin{tabular}{lrrp{0.38\linewidth}}
\toprule
Configuration & NanoMSMARCO & NanoNFCorpus & Notes \\
\midrule
SPLADE-v3 doc-doc ceiling & $\sim$0.64--0.65 & TODO & Full model on both sides. \\
\texttt{vocab\_transplant\_align} & $\sim$0.61 & TODO & Ettin-17M/68M to SPLADE-v3 vocab; alignment-only. \\
\texttt{splade\_shallow\_align} & $\sim$0.65 & TODO & Stable non-factorized shallow SPLADE. \\
\texttt{splade\_shallow\_align\_distill} & $\sim$0.65 & $\sim$0.31 & Good retrieval, unstable nnz/FLOPs. \\
\texttt{splade\_shallow\_factorized\_spaced\_align} & $\sim$0.693 & $\sim$0.333 & Strongest SPLADE-v3 factorized NanoBEIR run. \\
Lion-SP-1B doc-doc ceiling & 0.6510 & 0.3480 & Lion model on both sides. \\
\texttt{lion\_transplant\_align} & 0.6410 & 0.3211 & Ettin-150M transplanted to Lion 128K vocabulary. \\
\texttt{lion\_shallow\_factorized\_align}, 5 contiguous layers & 0.63--0.65 & 0.26--0.29 & Several runs; sensitive to factor dim and schedule. \\
\texttt{lion\_shallow\_factorized\_spaced\_align}, 5 spaced layers & TODO & TODO & Active run; fill with best checkpoint. \\
\bottomrule
\end{tabular}
\end{table}

\subsection{Ablations and Negative Results}

\begin{table}[t]
\caption{Ablations that clarify what is and is not responsible for success.}
\label{tab:ablations}
\centering
\begin{tabular}{p{0.25\linewidth}p{0.30\linewidth}p{0.34\linewidth}}
\toprule
Ablation & Outcome & Interpretation \\
\midrule
Random embedding init with donor vocabulary & Fails to match vocab transplant & Correct output dimensionality is not enough. \\
Frozen document MLM head plus projection & Fails & A pretrained lexical head alone is insufficient. \\
Shared-vocabulary BERT direct alignment & Underperforms & Sharing vocabulary without the right initialization/body is insufficient. \\
Ettin-68M vs Ettin-17M & Similar behavior & Capacity was not the first bottleneck in transplant runs. \\
More hard negatives ($n$-way 33) & No clear improvement & Ranking fine-tuning was not the binding constraint. \\
Cross-encoder MarginMSE & Competitive NanoBEIR, unstable nnz & Margins do not uniquely constrain sparse support. \\
Unfreezing factorized SPLADE lexical factors & Works but weaker & Moving the compressed lexical basis is risky. \\
Keeping factorized SPLADE head frozen & Strongest factorized result & Adapt the body, preserve the SVD lexical basis. \\
\bottomrule
\end{tabular}
\end{table}

\section{Discussion}

\subsection{Why Vocabulary Transplant Works, and Why It Breaks}

The transplant experiments show that sparse retrieval compatibility is not just
about output dimensionality. The embedding table is the interface between the
query backbone and the document vocabulary. Tokensurgeon-style kNN
interpolation gives the backbone a usable starting geometry; random
initialization does not.

The Lion transplant runs make the limitation clearer. Llama-3 tokenization uses
a large BPE vocabulary, while many BERT-family backbones use WordPiece. Exact
string overlap can be single-digit percent, so most donor tokens must be
interpolated. For robust modern backbones this can still work, but for
bert-base the transplanted embedding table appears too out-of-distribution: the
model can learn unigram-level masked language modeling, but does not recover
context-sensitive behavior quickly enough.

\subsection{Why Shallow Same-Checkpoint Models Are More Stable}

Shallow same-checkpoint models sidestep the vocabulary-transfer problem. Query
and document vectors are compatible by construction because they share the
tokenizer and lexical head. The remaining question is how much transformer depth
the query side needs.

The SPLADE-v3 factorized spaced result suggests that a query encoder does not
need a contiguous prefix of layers. A 3-layer body using one early, one middle,
and one final layer performed better than the contiguous factorized variants.
This motivates ongoing Lion experiments with 5 spaced layers.

\subsection{Sparse Support Stability}

Several runs achieved decent retrieval while producing undesirable sparse
support dynamics. This matters because sparse retrieval efficiency depends on
the number and mass of active dimensions, not just ranking metrics. The
cross-encoder MarginMSE experiment is the clearest example: score differences
can be matched with many query-vector shapes, so the model can vary nnz and mass
without being punished by the margin objective. Future losses should anchor
support or lexical mass directly if they use cross-encoder scores.

\subsection{Current Lion Status}

Lion-SP-1B has a much larger lexical matrix than SPLADE-v3. The full merged
model has approximately 1.236B parameters, and the tied lexical matrix alone is
approximately 262.7M parameters. Factorization is therefore not optional if the
goal is a query model in the 25--30\% parameter range. With
factorized dimension 256, a 5-layer Lion query model is roughly 337M
parameters, about 27\% of the full model. Increasing the factorization dimension
to 384 did not materially improve results in the current runs, so layer
selection is the more promising next axis.

\section{Future Experiment Template}

\begin{table}[t]
\caption{Template for upcoming experiments. Fill this table as new runs finish.}
\label{tab:future}
\centering
\begin{tabular}{p{0.25\linewidth}p{0.25\linewidth}p{0.18\linewidth}p{0.20\linewidth}}
\toprule
Experiment & Hypothesis & Expected metric & Result / decision \\
\midrule
Correct LR scheduler for gradient accumulation & Current cosine decay is too slow when scheduler steps only on optimizer updates. & Better late-stage NanoBEIR and full MS MARCO. & TODO \\
Lion 5 spaced layers [0,3,6,10,-1] & Later Lion layers help lexical/semantic alignment more than first-5 contiguous layers. & NanoMSMARCO $>$ contiguous 5-layer run. & Active; fill after 50k. \\
Lion full MS MARCO eval & NanoBEIR improvements transfer to full corpus. & NDCG@10, MRR@10. & TODO; requires Lion index. \\
SPLADE-v3 factorized repeat with fixed scheduler & Confirm best factorized result under corrected schedule. & NDCG@10 around or above 0.4438. & TODO \\
Cross-encoder MarginMSE with vector anchor & Stabilize nnz while using reranker margins. & Similar NDCG with lower nnz variance. & TODO \\
BPE-aware transplant & Improve Lion transplant for BERT/RoBERTa backbones. & Better MLM warmup and NanoBEIR. & TODO \\
Latency benchmark & Measure query-time savings directly. & ms/query and params. & TODO \\
\bottomrule
\end{tabular}
\end{table}

\section{Threats to Validity}

First, many results are from NanoBEIR development subsets. These are useful for
debugging but have high variance. Claims should rely on full MS MARCO dev
whenever possible.

Second, some training recipes changed during exploration. In particular, the
interaction between gradient accumulation and cosine learning-rate schedules
must be standardized before final claims.

Third, the strongest full result is currently for SPLADE-v3. Lion results are
promising but need full-index evaluation.

Fourth, parameter count is not the same as latency. The final paper should
include actual query encoding latency on at least one GPU and one CPU setting.

\section{Conclusion}

We studied how to train cheap query encoders for frozen sparse document models.
Vocabulary transplant showed that cross-vocabulary asymmetric SPLADE is
possible, but also revealed that embedding initialization quality is
load-bearing. Same-checkpoint shallow query encoders are more stable because
they preserve the sparse vocabulary by construction. The strongest completed
result combines shallow layer selection with ALBERT-style lexical factorization:
a 3-layer factorized SPLADE-v3 query encoder retains about 95\% of full
SPLADE-v3 MS MARCO dev quality. Current Lion experiments extend the same idea
to a larger 128K-vocabulary sparse document model; the main open question is
whether spaced layer selection and corrected training schedules can produce a
similarly strong full-corpus result.

\begin{thebibliography}{99}

\bibitem{Formal2021SPLADE}
T. Formal, B. Piwowarski, and S. Clinchant.
\newblock SPLADE: Sparse Lexical and Expansion Model for First Stage Ranking.
\newblock In \emph{SIGIR}, 2021.

\bibitem{Formal2022SPLADEv2}
T. Formal, C. Lassance, B. Piwowarski, and S. Clinchant.
\newblock From Distillation to Hard Negative Sampling: Making Sparse Neural IR Models More Effective.
\newblock In \emph{SIGIR}, 2022.

\bibitem{Lassance2024SPLADEv3}
C. Lassance, H. Dejean, T. Formal, and S. Clinchant.
\newblock SPLADE-v3: New baselines for SPLADE.
\newblock arXiv preprint, 2024.

\bibitem{Zeng2025Lion}
H. Zeng et al.
\newblock Scaling Sparse and Dense Retrieval in Decoder-Only LLMs.
\newblock TODO: verify final venue/version and complete author list. Model card:
\texttt{hzeng/Lion-SP-1B-llama3-marco-mntp}.

\bibitem{Lan2020ALBERT}
Z. Lan, M. Chen, S. Goodman, K. Gimpel, P. Sharma, and R. Soricut.
\newblock ALBERT: A Lite BERT for Self-supervised Learning of Language Representations.
\newblock In \emph{ICLR}, 2020.

\bibitem{Minixhofer2022WECHSEL}
B. Minixhofer, F. Paischer, and N. Rekabsaz.
\newblock WECHSEL: Effective initialization of subword embeddings for cross-lingual transfer of monolingual language models.
\newblock In \emph{NAACL}, 2022.

\bibitem{Bajaj2016MSMARCO}
P. Bajaj, D. Campos, N. Craswell, L. Deng, J. Gao, X. Liu, R. Majumder,
A. McNamara, B. Mitra, T. Nguyen, M. Rosenberg, X. Song, A. Stoica,
S. Tiwary, and T. Wang.
\newblock MS MARCO: A Human Generated MAchine Reading COmprehension Dataset.
\newblock arXiv:1611.09268, 2016.

\bibitem{Thakur2021BEIR}
N. Thakur, N. Reimers, A. Ruckl\'e, A. Srivastava, and I. Gurevych.
\newblock BEIR: A Heterogeneous Benchmark for Zero-shot Evaluation of Information Retrieval Models.
\newblock In \emph{NeurIPS Datasets and Benchmarks}, 2021.

\bibitem{Santhanam2022ColBERTv2}
K. Santhanam, O. Khattab, J. Saad-Falcon, C. Potts, and M. Zaharia.
\newblock ColBERTv2: Efficient and Effective Retrieval via Lightweight Late Interaction.
\newblock In \emph{NAACL}, 2022.

\bibitem{Zhuang2021TILDE}
S. Zhuang and G. Zuccon.
\newblock TILDE: Term Independent Likelihood moDEl for Passage Re-ranking.
\newblock In \emph{SIGIR}, 2021.

\bibitem{Lin2021UniCOIL}
J. Lin and X. Ma.
\newblock A Few Brief Notes on DeepImpact, COIL, and a Conceptual Framework for Information Retrieval Techniques.
\newblock arXiv:2106.14807, 2021.

\bibitem{Mallia2021DeepImpact}
A. Mallia, O. Khattab, T. Suel, and N. Tonellotto.
\newblock Learning Passage Impacts for Inverted Indexes.
\newblock In \emph{SIGIR}, 2021.

\bibitem{Shen2023LexMAE}
T. Shen et al.
\newblock LexMAE: Lexicon-Bottlenecked Pretraining for Large-Scale Retrieval.
\newblock In \emph{ICLR}, 2023.

\end{thebibliography}

\end{document}
